% !TEX program = pdflatex
\documentclass[11pt, a4paper]{article}

% ── Encoding & Fonts ──────────────────────────────────────────────
\usepackage[utf8]{inputenc}
\usepackage[T1]{fontenc}
\usepackage{lmodern}
\usepackage{microtype}

% ── Page layout ───────────────────────────────────────────────────
\usepackage[margin=2.4cm, top=2.8cm, bottom=2.8cm]{geometry}
\usepackage{setspace}
\onehalfspacing

% ── Math ──────────────────────────────────────────────────────────
\usepackage{amsmath, amssymb, amsthm, mathtools}
\usepackage{bm}

% ── Graphics & floats ────────────────────────────────────────────
\usepackage{graphicx}
\usepackage{xcolor}
\usepackage{booktabs}
\usepackage{float}
\usepackage{caption}
\usepackage{subcaption}

% ── References ────────────────────────────────────────────────────
\usepackage[numbers, sort&compress]{natbib}
\usepackage[colorlinks=true, linkcolor=blue!60!black,
            citecolor=green!50!black, urlcolor=blue!70!black]{hyperref}
\usepackage[capitalize, nameinlink]{cleveref}

% ── Theorem environments ─────────────────────────────────────────
\theoremstyle{plain}
\newtheorem{theorem}{Theorem}[section]
\newtheorem{proposition}[theorem]{Proposition}
\newtheorem{corollary}[theorem]{Corollary}
\newtheorem{lemma}[theorem]{Lemma}
\theoremstyle{definition}
\newtheorem{definition}[theorem]{Definition}
\theoremstyle{remark}
\newtheorem{remark}[theorem]{Remark}

% ── Custom commands ──────────────────────────────────────────────
\newcommand{\R}{\mathbb{R}}
\newcommand{\E}{\mathbb{E}}
\DeclareMathOperator*{\argmin}{arg\,min}
\DeclareMathOperator*{\argmax}{arg\,max}
\newcommand{\bx}{\bm{x}}
\newcommand{\by}{\bm{y}}
\newcommand{\bn}{\bm{n}}
\newcommand{\PSNR}{\mathrm{PSNR}}
\newcommand{\MSE}{\mathrm{MSE}}
\newcommand{\SSIM}{\mathrm{SSIM}}
\newcommand{\dB}{\,\mathrm{dB}}
\newcommand{\Jcal}{\mathcal{J}}
\newcommand{\Fcal}{\mathcal{F}}
\newcommand{\Ncal}{\mathcal{N}}
\newcommand{\Hcal}{\mathcal{H}}

\graphicspath{{figures/}}

% ══════════════════════════════════════════════════════════════════
\title{\bfseries Gaussian Splat Count vs.\ Reconstruction Quality:\\
       Rate-Distortion Analysis with Blind-Spot Cross-Validation}
\author{Supplementary Material --- Luxar Project}
\date{\today}

\begin{document}
\maketitle

% ──────────────────────────────────────────────────────────────────
\begin{abstract}
Gaussian splatting provides a compact representation of volumetric
microscopy data, but the choice of how many splats to use is typically
made by trial and error. We present a systematic analysis of the
rate-distortion trade-off for Gaussian splat representations across
12 microscopy datasets spanning confocal, spinning-disk confocal, and
light-sheet modalities. We introduce a principled model selection criterion
based on blind-spot cross-validation
\citep{batson2019noise2self}: by masking a fraction of voxels during
fitting and evaluating reconstruction on the held-out positions, we
determine the splat count at which the model transitions from fitting
signal to fitting noise. We provide a theoretical justification for
why the held-out PSNR peak (or plateau onset) identifies the optimal
splat count, and validate this across all datasets. Additionally, we
estimate the noise floor for each dataset using robust estimators
(Laplacian MAD and Haar wavelet MAD), providing a theoretical PSNR
ceiling that contextualises the rate-distortion curves. Our analysis
reveals that: (i) optimal splat counts range from 15K to 124K depending
on the dataset, (ii)~Gaussian splats exhibit clear overfitting on 10 of
12 datasets when given excessive capacity, and (iii)~the noise floor
and held-out peak together provide complementary stopping criteria for
automatic splat count selection.
\end{abstract}

\tableofcontents

% ══════════════════════════════════════════════════════════════════
\section{Introduction}
\label{sec:intro}

Gaussian splatting \citep{kerbl20233dgs} represents a volume
$V:\R^d\to\R_{\ge 0}$ as a weighted sum of $N$ oriented Gaussian kernels:
\begin{equation}
  \hat{V}(\bx) = \sum_{i=1}^{N} a_i \,
  \Ncal\!\left(\bx \mid \bm{\mu}_i, \bm{\Sigma}_i\right),
  \label{eq:gsplat-model}
\end{equation}
where $a_i > 0$ are amplitudes, $\bm{\mu}_i\in\R^d$ are centres, and
$\bm{\Sigma}_i \succ 0$ are covariance matrices parameterised via their
Cholesky factors. The model is fit by minimising a loss (typically $L^1$
or MSE) between $\hat{V}$ and the observed noisy volume
$Y = V + \epsilon$, where $\epsilon$ is pixel-independent noise.

A central practical question is: \emph{how many splats~$N$ should one use?}
Too few and the model underfits, missing genuine structure. Too many and
the model overfits, memorising noise that does not generalise. This
supplementary document addresses this question through three complementary
analyses:

\begin{enumerate}
  \item \textbf{Rate-distortion curves} (\cref{sec:rate-distortion}):
    PSNR and SSIM as a function of splat count, from 1K to 512K.
  \item \textbf{Blind-spot cross-validation} (\cref{sec:noise2self}):
    blind-spot cross-validation to identify the optimal splat count, with
    a formal theoretical justification.
  \item \textbf{Noise floor estimation} (\cref{sec:noise-floor}):
    robust estimation of the noise standard deviation and the theoretical
    PSNR ceiling.
\end{enumerate}

\subsection{Datasets}
\label{sec:datasets}

We evaluate on 12 datasets spanning 3 imaging modalities, 5 organisms,
and volume sizes from 4\,M to 100\,M voxels (\cref{tab:datasets}).

\begin{table}[H]
\centering
\caption{Dataset summary. All volumes are 3D float32 normalised to $[0,1]$.}
\label{tab:datasets}
\small
\begin{tabular}{@{}llllr@{}}
\toprule
Key & Modality & Structure & Shape & Voxels \\
\midrule
opencell\_map4\_ch0\cite{cho2022opencell} & Spinning-disk & Nuclei (Hoechst) & $51\times600\times600$ & 18.4\,M \\
opencell\_map4\_ch1 & Spinning-disk & Microtubules (MAP4) & $51\times600\times600$ & 18.4\,M \\
opencell\_lmnb1\_ch0 & Spinning-disk & Nuclei (Hoechst) & $51\times600\times600$ & 18.4\,M \\
opencell\_lmnb1\_ch1 & Spinning-disk & Nuclear lamina (LMNB1) & $51\times600\times600$ & 18.4\,M \\
kidney\_dapi & Confocal & Tissue nuclei (DAPI) & $16\times512\times512$ & 4.2\,M \\
kidney\_actin & Confocal & Actin filaments & $16\times512\times512$ & 4.2\,M \\
cells3d\_nuclei & Confocal & HeLa nuclei & $60\times256\times256$ & 3.9\,M \\
cells3d\_membrane & Confocal & HeLa membrane & $60\times256\times256$ & 3.9\,M \\
organoid\_ch0 & Confocal & Organoid & $236\times275\times271$ & 17.6\,M \\
celegans\_t100 & Confocal & Embryo nuclei & $41\times512\times512$ & 10.7\,M \\
tribolium & Light-sheet & Embryo (cropped) & $991\times104\times965$ & 99.5\,M \\
acto3d\_heart & Light-sheet & Heart nuclei (cropped) & cropped & $\le$100\,M \\
\bottomrule
\end{tabular}
\end{table}


% ══════════════════════════════════════════════════════════════════
\section{Rate-Distortion Analysis}
\label{sec:rate-distortion}

For each dataset, we fit Gaussian splats at 10 logarithmically spaced
counts: $N \in \{1\text{K}, 2\text{K}, 4\text{K}, 8\text{K}, 16\text{K},
32\text{K}, 64\text{K}, 128\text{K}, 256\text{K}, 512\text{K}\}$, using
non-progressive single-pass fitting with $L^1$ loss, 20\,000 maximum
iterations, early stopping (patience 500), and conservative culling
($99.9\%$ amplitude retention). For each count, we render the fitted
model back to the original volume shape and compute PSNR, SSIM
\citep{wang2004ssim}, MSE, and compression ratio.

The compression ratio is defined as
\begin{equation}
  \text{CR} = \frac{|\Omega|}{N \cdot (d + d(d+1)/2 + 1)},
\end{equation}
where $|\Omega|$ is the number of voxels and each 3D splat requires
$d + d(d+1)/2 + 1 = 10$ floats (3 centre + 6 Cholesky + 1 amplitude).

Results for all 12 datasets are shown in \cref{sec:figures}. Each
figure contains four panels: PSNR vs.\ splat count (with noise floor
overlay), SSIM vs.\ splat count, fitting time, and culling percentage.


% ══════════════════════════════════════════════════════════════════
\section{Blind-Spot Cross-Validation}
\label{sec:noise2self} % Label kept for cross-reference stability

\subsection{Motivation}
\label{sec:n2s-motivation}

The rate-distortion curves show PSNR monotonically increasing with splat
count when measured against the \emph{noisy} input. But this is misleading:
beyond some point, additional splats absorb noise rather than signal.
To detect this transition, we need a quality measure that evaluates
reconstruction of the \emph{underlying signal}, not the noise.

\subsection{Blind-Spot Cross-Validation}
\label{sec:n2s-method}

Inspired by the blind-spot framework of Batson \& Royer \citep{batson2019noise2self}, we
perform the following procedure for each splat count:

\begin{enumerate}
  \item \textbf{Mask}: Select 5\% of voxels uniformly at random
    (mask $M \subset \Omega$, $|M| = 0.05|\Omega|$).
  \item \textbf{Replace}: For each masked voxel $j \in M$, replace its
    value with the \emph{donut median} of its $3\times3\times3$
    neighbourhood excluding the centre (blind spot). This produces
    a modified volume $\tilde{Y}$.
  \item \textbf{Fit}: Fit Gaussian splats to $\tilde{Y}$ (the model
    never sees the original values at masked positions).
  \item \textbf{Evaluate}: Compute the \emph{held-out} MSE and PSNR
    at masked positions against the \emph{original} noisy values $Y_j$:
    \begin{equation}
      \MSE_{\text{held-out}} = \frac{1}{|M|}
        \sum_{j \in M} \left(\hat{V}(\bx_j) - Y_j\right)^2.
    \end{equation}
\end{enumerate}

\subsection{Theoretical Justification}
\label{sec:n2s-theory}

We now show that the splat count $N^*$ that maximises the held-out PSNR
is the one that best approximates the true denoised signal.

\begin{definition}[Observation model]
Let $V:\Omega \to \R$ be the true signal and $Y = V + \epsilon$ the
noisy observation, where $\epsilon_j \sim (0, \sigma^2)$ are independent,
zero-mean, with variance $\sigma^2$.
\end{definition}

\begin{definition}[$\Jcal$-invariance]
A function $f: \R^{|\Omega|} \to \R^{|\Omega|}$ is
\emph{$\Jcal$-invariant} with respect to a partition
$\Jcal = \{J_1, \ldots, J_K\}$ of $\Omega$ if, for each $k$, the
output $f(Y)_{J_k}$ does not depend on $Y_{J_k}$.
\end{definition}

Gaussian splatting is not exactly $\Jcal$-invariant because each splat
covers multiple voxels. However, our blind-spot masking procedure
\emph{enforces} $\Jcal$-invariance by construction: the model $\hat{V}$
is fitted to $\tilde{Y}$, where the values at masked positions $M$
have been replaced by their donut median. Thus, $\hat{V}(\bx_j)$ for
$j \in M$ does not depend on $Y_j$.

\begin{proposition}[Held-out loss decomposes into signal and noise terms]
\label{prop:held-out-decomposition}
For masked positions $M$ where the fitted model is $\Jcal$-invariant
(does not depend on $Y_j$), the expected held-out MSE decomposes as:
\begin{equation}
  \E\left[\MSE_{\text{held-out}}\right]
  = \underbrace{\frac{1}{|M|}\sum_{j\in M}
    \E\left[(\hat{V}(\bx_j) - V_j)^2\right]}_{\text{signal error}}
  + \sigma^2.
  \label{eq:held-out-decomposition}
\end{equation}
\end{proposition}

\begin{proof}
For each masked voxel $j \in M$:
\begin{align}
  \E\left[(\hat{V}(\bx_j) - Y_j)^2\right]
  &= \E\left[(\hat{V}(\bx_j) - V_j - \epsilon_j)^2\right] \\
  &= \E\left[(\hat{V}(\bx_j) - V_j)^2\right]
     - 2\E\left[(\hat{V}(\bx_j) - V_j)\epsilon_j\right]
     + \E[\epsilon_j^2].
\end{align}
Since $\hat{V}(\bx_j)$ is $\Jcal$-invariant with respect to $j$
(fitted on $\tilde{Y}$ where $Y_j$ was replaced), $\hat{V}(\bx_j)$
is independent of $\epsilon_j$. Furthermore, $V_j$ is deterministic
and $\E[\epsilon_j] = 0$. Therefore:
\begin{equation}
  \E\left[(\hat{V}(\bx_j) - V_j)\epsilon_j\right]
  = \E\left[\hat{V}(\bx_j) - V_j\right]\E[\epsilon_j] = 0.
\end{equation}
Summing over $M$ and dividing by $|M|$ yields
\cref{eq:held-out-decomposition}.
\end{proof}

\begin{corollary}[Optimal splat count minimises signal error]
\label{cor:optimal-N}
Since $\sigma^2$ is a constant independent of $N$, minimising the
held-out MSE (equivalently, maximising the held-out PSNR) over the
splat count $N$ is equivalent to minimising the signal reconstruction
error:
\begin{equation}
  N^* = \argmin_N\; \E\left[\MSE_{\text{held-out}}(N)\right]
      = \argmin_N\; \frac{1}{|M|}\sum_{j\in M}
        \E\left[(\hat{V}_N(\bx_j) - V_j)^2\right].
\end{equation}
\end{corollary}

\begin{remark}[Why the held-out PSNR peaks]
As $N$ increases from small values, the signal error decreases (more
splats capture more structure), so the held-out PSNR increases. Beyond
$N^*$, additional splats begin to fit patterns in the unmasked noise
that do not generalise to the masked positions, causing the signal error
at masked positions to increase (or plateau). The held-out PSNR thus
peaks at $N^*$ and declines (or flattens) thereafter.
\end{remark}

\begin{remark}[Connection to Batson \& Royer (2019)]
\Cref{prop:held-out-decomposition} is a special case of Proposition~1 in
\citet{batson2019noise2self}, applied to the specific setting where
$f = \hat{V}_N$ is a Gaussian splat model fitted to the blind-spot
volume $\tilde{Y}$. The key insight from the blind-spot framework is
that any $\Jcal$-invariant estimator's self-supervised loss (evaluated
against the noisy input) equals the supervised loss (evaluated against
the clean signal) plus a noise constant. This allows model selection
without access to clean data.
\end{remark}

\begin{remark}[Practical detection of $N^*$]
\label{rem:detection}
In practice, we detect $N^*$ using a hybrid criterion:
\begin{itemize}
  \item If the held-out PSNR curve has a \emph{clear peak} (the average
    of points before and after the argmax are both $\ge 0.1\dB$ below
    the peak), we use the argmax directly.
  \item Otherwise (plateau or monotonic curve), we report the first
    point within $0.3\dB$ of the maximum --- the onset of diminishing
    returns.
\end{itemize}
This hybrid approach is robust across all curve shapes observed in our
12 datasets.
\end{remark}

\subsection{Results}
\label{sec:n2s-results}

The cross-validation analysis reveals that \textbf{10 of 12 datasets exhibit
overfitting} when given excessive splat capacity. The held-out PSNR
peaks and then declines, confirming that additional splats memorise
noise. Only the two light-sheet datasets (Tribolium, Acto3D) show
monotonically increasing held-out PSNR, consistent with their very low
noise levels ($\sigma < 0.001$).

The optimal splat count $N^*$ varies by an order of magnitude across
datasets, from 15K (C.~elegans, sparse signal) to 124K (kidney actin,
dense tissue). See \cref{tab:n2s-summary} and the per-dataset figures.

\begin{table}[H]
\centering
\caption{Cross-validation model selection summary.}
\label{tab:n2s-summary}
\small
\begin{tabular}{@{}lrrrl@{}}
\toprule
Dataset & Peak PSNR & $N^*$ (splats) & Drop & Type \\
\midrule
opencell\_map4\_ch0 & 28.0\dB & 32K & 0.3\dB & Peak \\
opencell\_map4\_ch1 & 24.9\dB & 62K & 0.2\dB & Peak \\
opencell\_lmnb1\_ch0 & 25.3\dB & 31K & 0.3\dB & Peak \\
opencell\_lmnb1\_ch1 & 33.4\dB & 31K & 0.2\dB & Plateau \\
kidney\_dapi & 33.6\dB & 63K & 2.2\dB & Peak \\
kidney\_actin & 30.8\dB & 124K & 2.5\dB & Peak \\
cells3d\_nuclei & 34.7\dB & 63K & 0.3\dB & Peak \\
cells3d\_membrane & 45.5\dB & 63K & 1.9\dB & Peak \\
organoid\_ch0 & 37.7\dB & 28K & 0.6\dB & Peak \\
celegans\_t100 & 39.3\dB & 15K & 0.3\dB & Peak \\
tribolium & 42.8\dB & 31K & ---  & Plateau \\
acto3d\_heart & 40.5\dB & 400K & --- & Plateau \\
\bottomrule
\end{tabular}
\end{table}


% ══════════════════════════════════════════════════════════════════
\section{Noise Floor Estimation}
\label{sec:noise-floor}

\subsection{Method}

We estimate the noise standard deviation $\sigma$ using an ensemble of
three robust estimators:

\paragraph{Laplacian MAD} (after \citet{immerkaer1996fast}):
Apply the discrete Laplacian $L = \nabla^2 Y$, then estimate $\sigma$
via the median absolute deviation (MAD):
\begin{equation}
  \hat{\sigma}_{\text{Lap}} = \frac{\text{MAD}(L)}{0.6745 \cdot \sqrt{K}},
  \qquad K = 2d(2d+1),
\end{equation}
where $K$ is the sum of squared Laplacian kernel weights and $d$ is
the spatial dimension.

\paragraph{Haar finite-difference MAD} (after \citet{donoho1994ideal}):
For each $z$-slice, compute the 2D diagonal detail
$D_{y,x} = Y_{y,x} - Y_{y,x+1} - Y_{y+1,x} + Y_{y+1,x+1}$
(algebraically equivalent to the Haar wavelet HH subband). Then:
\begin{equation}
  \hat{\sigma}_{\text{Haar}} = \frac{\text{median}(|D|)}{0.6745 \cdot 2}.
\end{equation}

\paragraph{Background region MAD}: Identify voxels below the 10th
intensity percentile and compute $\hat{\sigma}_{\text{bg}} =
\text{MAD}(\text{background}) / 0.6745$.

The \textbf{ensemble estimate} is $\hat{\sigma} =
\text{median}(\hat{\sigma}_{\text{Lap}}, \hat{\sigma}_{\text{Haar}},
\hat{\sigma}_{\text{bg}})$, and the theoretical PSNR ceiling is:
\begin{equation}
  \PSNR_{\max} = -20\log_{10}\hat{\sigma}
  \qquad \text{(for data normalised to $[0,1]$)}.
  \label{eq:psnr-max}
\end{equation}

\subsection{Results}

\Cref{tab:noise-floor} shows the noise floor estimates. The Laplacian
and Haar methods agree within 15\% for all datasets; the background
method consistently underestimates (expected, as the darkest voxels have
less variance). The noise floor ranges from 30\dB\ (OpenCell ch1,
spinning-disk confocal) to 61\dB\ (Tribolium, light-sheet), reflecting
the superior noise characteristics of light-sheet over confocal imaging.

\begin{table}[H]
\centering
\caption{Noise floor estimates per dataset.}
\label{tab:noise-floor}
\small
\begin{tabular}{@{}lrrr@{}}
\toprule
Dataset & $\hat{\sigma}$ & $\PSNR_{\max}$ (dB) & Modality \\
\midrule
opencell\_map4\_ch0 & 0.0234 & 32.6 & Spinning-disk \\
opencell\_map4\_ch1 & 0.0300 & 30.4 & Spinning-disk \\
opencell\_lmnb1\_ch0 & 0.0319 & 30.0 & Spinning-disk \\
opencell\_lmnb1\_ch1 & 0.0075 & 42.6 & Spinning-disk \\
kidney\_dapi & 0.0086 & 41.3 & Confocal \\
kidney\_actin & 0.0067 & 43.5 & Confocal \\
cells3d\_nuclei & 0.0133 & 37.6 & Confocal \\
cells3d\_membrane & 0.0021 & 53.5 & Confocal \\
organoid\_ch0 & 0.0019 & 54.4 & Confocal \\
celegans\_t100 & 0.0092 & 40.7 & Confocal \\
tribolium & 0.0009 & 61.0 & Light-sheet \\
acto3d\_heart & $<$0.001 & $>$60 & Light-sheet \\
\bottomrule
\end{tabular}
\end{table}


% ══════════════════════════════════════════════════════════════════
\section{Discussion}
\label{sec:discussion}

\subsection{Key Findings}

\paragraph{1. Gaussian splats overfit on most microscopy datasets.}
With conservative culling (99.9\% amplitude retention), 10 of 12 datasets
show a clear decline in held-out PSNR beyond the optimal count $N^*$.
The two exceptions are light-sheet datasets (Tribolium, Acto3D) with
very low noise ($\sigma < 0.001$), where the model remains signal-limited
across all tested counts. This demonstrates that the implicit smoothness
of Gaussian basis functions does not fully prevent overfitting ---
contrary to what one might expect from a representation that cannot
represent pixel-independent noise.

\paragraph{2. The optimal splat count is dataset-dependent.}
$N^*$ ranges from 15K (C.~elegans) to 124K (kidney actin), spanning
an order of magnitude. The primary determinants are:
\begin{itemize}
  \item \emph{Signal complexity}: Dense tissue (kidney) requires more
    splats than sparse nuclei (C.~elegans).
  \item \emph{Noise level}: Higher noise (OpenCell, $\sigma \approx 0.03$)
    means fewer ``useful'' splats before overfitting.
  \item \emph{Volume geometry}: Thin volumes (kidney, 16 $z$-slices)
    overfit earlier than isotropic volumes.
\end{itemize}

\paragraph{3. The noise floor contextualises the rate-distortion curves.}
Datasets far below their noise floor (OpenCell: 30\dB\ vs.\ 33\dB\
ceiling) still have recoverable signal. Datasets near their noise floor
(C.~elegans: 39\dB\ vs.\ 41\dB\ ceiling) are nearly saturated. The
noise floor line in each figure provides an at-a-glance assessment of
how much room for improvement remains.

\paragraph{4. The CV held-out peak is a principled stopping criterion.}
\Cref{prop:held-out-decomposition} shows that the held-out PSNR directly
estimates signal reconstruction quality (up to a noise constant). Unlike
heuristics (e.g., ``use $N = |\Omega|/500$ splats''), the held-out peak
is adaptive to the specific dataset's noise level and signal complexity.

\subsection{Limitations}

\paragraph{Noise independence assumption.}
The blind-spot framework requires pixel-independent noise. Preprocessing
steps such as percentile clipping (applied to OpenCell and C.~elegans
datasets) can violate this assumption by introducing spatial correlations
at clipped extremes. This may partially explain the relatively flat
held-out curves for some datasets. For the most rigorous results, the
analysis should be applied to truly raw data.

\paragraph{Discrete splat count sampling.}
We evaluate at 10 geometrically spaced counts. The true $N^*$ lies
between sampled points; finer sampling near the detected peak would
improve precision.

% ══════════════════════════════════════════════════════════════════
\section{Conclusion}
\label{sec:conclusion}

We have presented a comprehensive rate-distortion analysis of Gaussian
splat representations for 3D microscopy, covering 12 datasets across
three imaging modalities. The blind-spot cross-validation (inspired by Batson \& Royer, 2019)
provides a principled, fully automatic criterion for selecting the optimal
splat count --- one that requires no clean reference data, no prior on the
noise level, and no manual parameter tuning. The theoretical justification
(\cref{prop:held-out-decomposition}) shows that maximising the held-out
PSNR is equivalent to minimising the signal reconstruction error, making
it the natural objective for model selection.

Our noise floor analysis complements this by establishing the theoretical
PSNR ceiling for each dataset, contextualising where the reconstruction
stands relative to the best achievable quality. Together, these tools
enable practitioners to make informed decisions about splat budgets for
their specific data.

% ══════════════════════════════════════════════════════════════════
% FIGURES — one page per dataset with 3 panels
% ══════════════════════════════════════════════════════════════════
\clearpage
\appendix
\section{Per-Dataset Figures}
\label{sec:figures}

% Helper macro for per-dataset figure pages
\newcommand{\datasetfigs}[2]{%
  \clearpage
  \subsection{#2}
  \begin{figure}[H]
    \centering
    \includegraphics[width=\textwidth]{#1_fig_quality_curves.pdf}
    \caption{Rate-distortion curves for #2.}
  \end{figure}
  \begin{figure}[H]
    \centering
    \includegraphics[width=\textwidth]{#1_fig_slice_montage.pdf}
    \caption{Reconstruction slice montage for #2.}
  \end{figure}
  \begin{figure}[H]
    \centering
    \includegraphics[width=\textwidth]{#1_fig_noise2self.pdf}
    \caption{Cross-validation analysis for #2.}
  \end{figure}
}

\datasetfigs{opencell_map4_ch0}{MAP4 (Hoechst) --- Nuclei}
\datasetfigs{opencell_map4_ch1}{MAP4 (GFP) --- Microtubules}
\datasetfigs{opencell_lmnb1_ch0}{LMNB1 (Hoechst) --- Nuclei}
\datasetfigs{opencell_lmnb1_ch1}{LMNB1 (GFP) --- Nuclear Lamina}
\datasetfigs{kidney_dapi}{Mouse Kidney --- DAPI (Nuclei)}
\datasetfigs{kidney_actin}{Mouse Kidney --- Phalloidin (Actin)}
\datasetfigs{cells3d_nuclei}{HeLa Cells --- Nuclei}
\datasetfigs{cells3d_membrane}{HeLa Cells --- Membrane}
\datasetfigs{organoid_ch0}{Organoid --- Channel 0}
\datasetfigs{celegans_t100}{C.\ elegans Embryo --- $t=100$}
\datasetfigs{tribolium}{Tribolium Embryo (Light-Sheet)}
\datasetfigs{acto3d_heart_nuclei}{Mouse Heart --- Nuclei (Light-Sheet)}

% ══════════════════════════════════════════════════════════════════
\bibliographystyle{unsrtnat}
\bibliography{references}

\end{document}
